\begin{abstract}
QSAR models for Ames mutagenicity are increasingly relied upon in regulatory assessments 
under EU REACH, K-REACH, and ICH M7, yet their reported performance varies widely across 
evaluation protocols, making it difficult to assess their true predictive 
reliability.\cite{Honma_2019,Furuhama_2023,Wu_MoleculeNet_2018} This study systematically 
evaluated 375 model configurations---seven algorithms, five feature representations, and 
three preprocessing scenarios---across three primary and two exploratory genotoxicity 
datasets: Ames ($n = 13{,}560$), in-vitro chromosome aberration ($n = 1{,}619$), in-vivo 
micronucleus ($n = 2{,}153$), and two negative-sampled variants. Using a single locked 
configuration (XGBoost, compact features, raw SMILES) in an evaluation cascade of 
increasing rigor, Ames MCC fell from 0.670 (random-split CV) to 0.624 (scaffold-split CV) 
to 0.234 and 0.122 (drug$\to$industrial and industrial$\to$drug leave-domain-out, LDO). 
The split-method effect ($\Delta = 0.046$) was nearly an order of magnitude smaller than 
the data source heterogeneity effect ($\Delta = 0.390$): a naive source classifier reached 
MCC~$= 0.608$ without any structural information, reflecting a 19.5-fold prevalence gap 
between drug-sourced (58.3\%) and industrial-sourced (3.0\%) compounds. Threshold 
decomposition showed that prior shift accounts for only part of this gap, with oracle MCC 
plateauing at 0.27--0.36. Algorithm robustness under source shift was uneven: Random Forest
maintained MCC~$= 0.325$ in the hardest direction where XGBoost collapsed to 0.122. 
In-vivo micronucleus models showed ranking ability (ROC-AUC~$= 0.860$, PR-AUC~$= 0.237$) 
but unreliable classification (MCC 95\% CI crossing zero). These findings support an 
evaluation cascade with source-awareness as a reporting standard for Ames mutagenicity QSAR.
\end{abstract}
